\chapter{Conclusion}
\label{ch:conclusion}

This thesis set out to address a specific limitation of the software project management platforms in wide use today: the absence of an independently verifiable, tamper-evident record of a project's history, stemming from their reliance on a single centralized provider to store and control that history (\Cref{sec:motivation}). Building on the architectural pattern validated by \textcite{adel2023} for decentralized data management in construction (\Cref{sec:adel}), this thesis adapted that pattern to the software project management domain and implemented it as a working prototype.

The resulting system models projects and tasks as chaincode-governed state on a permissioned Hyperledger Fabric network, exposing operations for project creation, role-based membership, task assignment, and status transitions, all recorded in Fabric's immutable transaction history rather than a provider-controlled log (\Cref{sec:chaincode-design}). File attachments are stored on IPFS and anchored to the ledger through their content identifiers, avoiding the cost of storing binary data on-chain while preserving a verifiable link between the recorded reference and the exact file content (\Cref{sec:ipfs-integration}). A backend API and browser-based frontend expose this functionality through a task-board workflow familiar to users of tools such as Trello or Jira, while making the system's central guarantee — the immutable audit trail — directly visible through a dedicated history view (\Cref{sec:backend,sec:frontend}).

The evaluation in \Cref{ch:results} confirmed that this design works as intended: project and task operations behave correctly and enforce membership constraints at the chaincode level, the audit trail correctly preserves every past state of a task, and file attachments are correctly anchored to and retrievable from IPFS via the CID recorded on-chain. Taken together, these results show that the Fabric-plus-IPFS pattern, previously validated only in a construction-industry context, transfers successfully to software project management without requiring a fundamentally different architecture — only a different data model and user-facing workflow.

\section{Limitations}
\label{sec:limitations}

The prototype was evaluated on a local, two-organization Fabric network rather than a geographically distributed or production-scale deployment, and this thesis does not report transaction throughput or latency under realistic multi-user load; the results in \Cref{ch:results} therefore demonstrate functional correctness rather than production readiness. The chaincode's data model, while sufficient to demonstrate the core architecture, is intentionally minimal: it does not yet cover features present in mature platforms such as sub-tasks, comments, notifications, or fine-grained per-field access control. Identity management in the prototype relies directly on Fabric's own X.509-based MSP identities rather than a separate user-facing authentication layer, which is adequate for demonstrating the architecture but would need to be extended for deployment to non-technical end users. Finally, the determinism constraint discussed in \Cref{sec:chaincode-design} is a reminder that this class of system carries development costs — around correctness patterns specific to multi-peer execution — that centralized platforms do not.

\section{Future Work}
\label{sec:future-work}

Several directions follow naturally from the limitations above. Extending the chaincode's data model to support sub-tasks, comments, and richer notification mechanisms would close the functional gap with mature centralized platforms while preserving the same on-chain audit guarantees. Evaluating the system on a geographically distributed network, with a realistic number of organizations and peers, would allow throughput and latency to be measured under conditions closer to a production deployment. Integrating a user-facing identity and authentication layer on top of Fabric's MSP would make the system usable by non-technical teams without requiring them to manage X.509 certificates directly. Finally, since the core architectural pattern evaluated in this thesis is not specific to software projects, applying the same chaincode-plus-IPFS approach to other collaborative domains with similar auditability requirements would help establish how broadly this pattern generalizes beyond the two domains — construction and software project management — in which it has now been demonstrated.
